
\section{Related Work}

\paragraph{Empirical analyses of on-policy distillation.}
Recent studies examine OPD through teacher--student compatibility, training data, and parameter updates. \citet{li2026rethinking} identify compatible thinking patterns and additional teacher capabilities as conditions for successful distillation, and trace alignment on shared high-probability tokens. \citet{fu2026rethinking} show broad state coverage from few queries and slowing supervision absorption, while \citet{hou2026matters} connect data efficiency to the longer reasoning trajectories induced by hard examples. In parameter space, \citet{shen2026geometry} identify early subspace locking and test constrained continuation; \citet{yu2026dense} characterize coordinate-sparse, spectrally concentrated updates distributed across layers. Further analyses examine sampled-token supervision failures \citep{fu2026revisiting} and associate training plateaus with weakening gradient signals \citep{zhao2026vanishing}. Test-time scaling reveals that OPD gains depend on the sampling budget \citep{ge2026testtime}, while controlled token corrections measure where teacher guidance can improve answers \citep{li2026demystifying}. Our study complements these findings by tracing teacher-induced gradient strength and distribution alongside hidden-representation displacement and the validation gains of actual optimizer steps. We further evaluate the performance and computational cost of three training-efficiency interventions.

\paragraph{Efficient OPD training.}
Existing methods improve supervision utilization through token selection, shorter rollouts, and response reuse. TIP uses student uncertainty and teacher--student disagreement \citep{xu2026tip}; rock-token interventions assess the reasoning contribution of persistently high-loss positions \citep{jiang2026rock}, and gradient-utility analysis evaluates alignment with task-success gradients \citep{armandpour2026unmasking}. Related approaches account for sampled-gradient estimation reliability \citep{sheng2026ier} or derive token weights from achievable loss reduction under a local KL budget \citep{sadagopan2026utility}. Prefix distillation shortens generation and supervises initial response segments \citep{zhang2026prefix}; KL Agreement Trap Termination cuts off rollouts that enter persistent low-KL degenerate agreement \citep{xin2026kat}. DistiLLM uses adaptive off-policy response reuse \citep{ko2024distillm}, while AsyncOPD decouples rollout generation from learner updates and addresses stale-policy supervision \citep{kang2026async}. Our experiments examine complementary decisions during continued OPD training: when to stop full-parameter updates, whether a representation-guided compact correction can support further learning, and how much additional fitting a cached response provides. 